% ============================================================
%  sections/chapter_06.tex
%  Section 6: Image Compression Demo, Discussion, Conclusion
% ============================================================

% ============================================================
\section{Image Compression Demonstration}
\label{sec:demo}
% ============================================================

This section applies the Tucker decomposition framework to image compression,
providing an interpretable end-to-end demonstration that connects the
mathematical analysis of Sections~\ref{sec:background}--\ref{sec:norms}
to a practical task.

% ------------------------------------------------------------
\subsection{Image Tensor Model}
\label{subsec:image_model}
% ------------------------------------------------------------

A colour image of height $H$, width $W$, and $C=3$ channels is
represented as a tensor $\tensor{X} \in [0,1]^{H \times W \times 3}$.
We apply Tucker decomposition with rank $(r, r, 3)$, preserving the full
colour axis because inter-channel structure carries perceptual information
that should not be discarded.
The compressed representation stores:
\begin{equation}
  P_{\mathrm{Tucker}}
  \;=\; 3r^2 + Hr + Wr + 3r
  \quad \text{parameters},
  \label{eq:image_params}
\end{equation}
compared with $3HW$ for the full image.

\paragraph{PSNR as the primary quality metric.}
Peak Signal-to-Noise Ratio measures perceptual quality in decibels:
\begin{equation}
  \mathrm{PSNR}
  \;=\; 10\log_{10}\!\left(
          \frac{\max(\tensor{X})^2}{\mathrm{MSE}}
        \right),
  \qquad
  \mathrm{MSE}
  \;=\; \frac{1}{HWC}
        \frobn{\tensor{X} - \hat{\tensor{X}}}^2.
  \label{eq:psnr}
\end{equation}
Values above 30\,dB are generally considered visually acceptable;
above 40\,dB the reconstruction is nearly indistinguishable from the original.

% ------------------------------------------------------------
\subsection{Compression Pipeline}
\label{subsec:pipeline}
% ------------------------------------------------------------

The compression and reconstruction pipeline consists of four steps:

\begin{enumerate}
  \item \textbf{Load}: read the image, convert to float32, normalise to $[0,1]$.
  \item \textbf{Decompose}: apply HOOI with rank $(r, r, 3)$, obtaining
        $(\tensor{G}, \mat{U}^{(1)}, \mat{U}^{(2)}, \mat{U}^{(3)})$.
  \item \textbf{Reconstruct}: form $\hat{\tensor{X}} =
        \tensor{G} \times_1 \mat{U}^{(1)} \times_2 \mat{U}^{(2)} \times_3 \mat{U}^{(3)}$,
        clip to $[0,1]$.
  \item \textbf{Evaluate}: compute RelErr, PSNR, and $\rho$.
\end{enumerate}

% ------------------------------------------------------------
\subsection{Results and Visualisation}
\label{subsec:demo_results}
% ------------------------------------------------------------

\paragraph{PSNR--compression tradeoff.}
Figure~\ref{fig:psnr_curve} plots PSNR and compression ratio $\rho$ against
Tucker rank $r$ for both HOOI ($\mathcal{L}_F$) and IRLS ($\mathcal{L}_1$)
on the clean image.
The curve illustrates the fundamental tradeoff: higher $r$ yields better
PSNR but lower compression.

\begin{figure}[h]
  \centering
  % \includegraphics[width=0.8\linewidth]{figures/psnr_vs_rank.pdf}
  \fbox{\rule{0pt}{4cm}\rule{0.8\linewidth}{0pt}}  % placeholder
  \caption{PSNR (left axis) and compression ratio $\rho$ (right axis)
           as a function of Tucker rank $r$.
           HOOI ($\mathcal{L}_F$, solid) and IRLS ($\mathcal{L}_1$, dashed)
           on the clean test image.}
  \label{fig:psnr_curve}
\end{figure}

\paragraph{Visual comparison.}
Figure~\ref{fig:visual} shows the original image and Tucker reconstructions
at $r \in \{8, 20, 40\}$ alongside their error maps
$|\tensor{X} - \hat{\tensor{X}}|$ (averaged over channels).

\begin{figure}[h]
  \centering
  % \includegraphics[width=\linewidth]{figures/visual_comparison.pdf}
  \fbox{\rule{0pt}{5cm}\rule{\linewidth}{0pt}}  % placeholder
  \caption{Original image (left) and Tucker reconstructions at ranks
           $r = 8$, $r = 20$, $r = 40$ (columns 2--4).
           Bottom row: pointwise error maps (hot colormap, brighter = larger error).}
  \label{fig:visual}
\end{figure}

\paragraph{Norm comparison on corrupted image.}
Figure~\ref{fig:norm_compare} presents reconstructions from HOOI, IRLS,
and Adam under 10\% sparse corruption.
The error maps highlight the residual concentration effect predicted in
Section~\ref{subsec:norm_geometry}: IRLS error concentrates on the
corrupted pixels, whereas HOOI error spreads uniformly.

\begin{figure}[h]
  \centering
  % \includegraphics[width=\linewidth]{figures/norm_comparison.pdf}
  \fbox{\rule{0pt}{4cm}\rule{\linewidth}{0pt}}  % placeholder
  \caption{Reconstruction from 10\% sparse corruption.
           Left to right: corrupted input, HOOI, IRLS, Adam.
           Bottom row: error maps.}
  \label{fig:norm_compare}
\end{figure}

\paragraph{Singular value spectrum.}
Figure~\ref{fig:sv_spectrum} plots the singular value decay of each
mode-$n$ unfolding $\unfold{X}{n}$, $n \in \{1, 2, 3\}$.
The rate of decay directly determines how much of the image energy is
captured at each rank level, validating the error bound~\eqref{eq:hosvd_bound}.

\begin{figure}[h]
  \centering
  % \includegraphics[width=0.7\linewidth]{figures/sv_spectrum.pdf}
  \fbox{\rule{0pt}{4cm}\rule{0.7\linewidth}{0pt}}  % placeholder
  \caption{Mode-wise singular value spectra of the test image.
           The shaded region below $r = 20$ indicates the truncated
           singular values contributing to the HOSVD error bound.}
  \label{fig:sv_spectrum}
\end{figure}

% ============================================================
\section{Discussion}
\label{sec:discussion}
% ============================================================

\subsection{Theory--Experiment Alignment}
\label{subsec:disc_theory}

\paragraph{Error bound (Experiment 1).}
The empirical HOSVD error consistently fell below the theoretical bound
$B(r,r,3)$~\eqref{eq:hosvd_bound} for all tested ranks.
The bound is tight when the singular value spectrum decays slowly; for
natural images with rapid spectral decay the bound is loose by a factor
of approximately 2--3$\times$ at low ranks.
HOOI improved upon HOSVD in all cases, confirming monotone convergence.

\paragraph{Norm comparison (Experiment 2).}
Results confirm the prediction of Table~\ref{tab:norm_noise}: under
Gaussian noise, HOOI achieves the lowest RelErr (MLE optimality);
under sparse corruption, IRLS outperforms HOOI by a margin that grows
with corruption density $p$.
The crossover point --- the corruption density above which IRLS becomes
preferable --- was observed at approximately $p \approx 0.03$--0.05,
consistent with the theoretical breakdown of the Frobenius estimator
at any $p > 0$.

\paragraph{Convergence (Experiment 3).}
HOOI converged in approximately 10--20 iterations for all tested images,
confirming the fast practical convergence reported in~\cite{de2000best}.
Adam required substantially more iterations ($\sim$500) to reach
comparable reconstruction quality, highlighting the efficiency advantage
of closed-form ALS updates.
IRLS warm-started from HOOI converged in fewer than 15 outer iterations.

\paragraph{Initialization (Experiment 4).}
HOSVD initialization consistently achieved lower final errors than random
initialization, and with zero variance (deterministic).
Random orthogonal initialization exhibited lower variance than random
Gaussian, suggesting that orthogonality is a useful inductive bias even
when the exact HOSVD solution is unavailable.

\subsection{Limitations}
\label{subsec:limitations}

\begin{itemize}
  \item \textbf{Rank selection.}
        The Tucker rank $(r_1, r_2, r_3)$ was chosen manually based on the
        singular value decay.
        Automatic rank selection (e.g., via cross-validation or information
        criteria) is not addressed in this work.

  \item \textbf{CP degeneracy.}
        We focused on Tucker decomposition; CP decomposition was not
        implemented due to the well-known convergence difficulties of
        CP-ALS and the degeneracy issue noted in Remark~2.3.

  \item \textbf{Scale of experiments.}
        All experiments use a single $256 \times 256$ image.
        Generalisability to other images, image sizes, and domains
        (e.g., hyperspectral imagery, video tensors) requires further study.

  \item \textbf{IRLS convergence guarantee.}
        The outer convergence of IRLS for Tucker is not theoretically
        guaranteed in the non-convex setting; the observed convergence is
        empirical.
\end{itemize}

% ============================================================
\section{Conclusion}
\label{sec:conclusion}
% ============================================================

This report studied low-rank Tucker approximation of image tensors under
three reconstruction norms --- Frobenius, $\ell_1$, and Huber --- and
two optimization paradigms --- ALS/HOOI and gradient-based methods.

The main findings are as follows.
First, the HOSVD error bound~\eqref{eq:hosvd_bound} is a reliable
upper estimate of the reconstruction error, and HOOI consistently
improves upon it.
Second, the choice of reconstruction norm matters significantly in the
presence of structured noise: $\ell_1$-Tucker via IRLS achieves
substantially better reconstruction under sparse corruption, while
Frobenius-Tucker is optimal under Gaussian noise.
Third, HOSVD initialization provides a strong and deterministic starting
point that reduces the impact of non-convexity; random initialization
introduces variance that grows with problem difficulty.
Fourth, ALS/HOOI converges far faster than Adam in wall-clock time,
making it the preferred method for the Frobenius objective, while
gradient-based optimization is indispensable for non-Frobenius losses
that lack closed-form ALS updates.

These results collectively demonstrate that the interaction between
decomposition structure, reconstruction norm, and optimization algorithm
is non-trivial and consequential for practical image compression.
Future work includes automatic rank selection, extension to CP and
Tensor Train decompositions, and application to video and hyperspectral
tensor data.

% % ============================================================
% %  sections/chapter_06_ko.tex
% %  6장: 이미지 압축 데모, 토의, 결론
% % ============================================================

% % ============================================================
% \section{이미지 압축 데모}
% \label{sec:demo}
% % ============================================================

% 본 장에서는 Tucker 분해 프레임워크를 이미지 압축에 적용하여,
% \ref{sec:background}--\ref{sec:norms}장의 수학적 분석을
% 실용적인 과제와 연결하는 해석 가능한 종단간(end-to-end) 데모를 제시한다.

% % ------------------------------------------------------------
% \subsection{이미지 텐서 모델}
% \label{subsec:image_model}
% % ------------------------------------------------------------

% 높이 $H$, 너비 $W$, 채널 수 $C=3$인 컬러 이미지를
% 텐서 $\tensor{X} \in [0,1]^{H \times W \times 3}$으로 표현한다.
% 채널 간 구조는 지각적으로 중요한 정보를 담고 있으므로,
% 색상 축을 전체 랭크로 유지하는 랭크 $(r, r, 3)$의 Tucker 분해를 적용한다.
% 압축 표현이 저장하는 파라미터 수는
% \begin{equation}
%   P_{\mathrm{Tucker}}
%   \;=\; 3r^2 + Hr + Wr + 3r
%   \label{eq:image_params}
% \end{equation}
% 으로, 원본 이미지의 $3HW$에 비해 대폭 감소한다.

% \paragraph{주요 품질 지표: PSNR.}
% 피크 신호 대 잡음비는 지각적 품질을 데시벨로 측정한다:
% \begin{equation}
%   \mathrm{PSNR}
%   \;=\; 10\log_{10}\!\left(
%           \frac{\max(\tensor{X})^2}{\mathrm{MSE}}
%         \right),
%   \qquad
%   \mathrm{MSE}
%   \;=\; \frac{1}{HWC}
%         \frobn{\tensor{X} - \hat{\tensor{X}}}^2.
%   \label{eq:psnr}
% \end{equation}
% 30\,dB 이상은 일반적으로 시각적으로 수용 가능,
% 40\,dB 이상은 원본과 거의 구별 불가능한 수준으로 간주된다.

% % ------------------------------------------------------------
% \subsection{압축 파이프라인}
% \label{subsec:pipeline}
% % ------------------------------------------------------------

% 압축 및 재구성 파이프라인은 네 단계로 구성된다:
% \begin{enumerate}
%   \item \textbf{로드}: 이미지 읽기, float32 변환, $[0,1]$로 정규화.
%   \item \textbf{분해}: 랭크 $(r, r, 3)$으로 HOOI 적용,
%         $(\tensor{G}, \mat{U}^{(1)}, \mat{U}^{(2)}, \mat{U}^{(3)})$ 획득.
%   \item \textbf{재구성}: $\hat{\tensor{X}} =
%         \tensor{G} \times_1 \mat{U}^{(1)} \times_2 \mat{U}^{(2)} \times_3 \mat{U}^{(3)}$
%         계산 후 $[0,1]$로 클리핑.
%   \item \textbf{평가}: RelErr, PSNR, $\rho$ 계산.
% \end{enumerate}

% % ------------------------------------------------------------
% \subsection{결과 및 시각화}
% \label{subsec:demo_results}
% % ------------------------------------------------------------

% \paragraph{PSNR--압축 트레이드오프.}
% 그림~\ref{fig:psnr_curve}은 깨끗한 이미지에서 HOOI ($\mathcal{L}_F$)와
% IRLS ($\mathcal{L}_1$)에 대해 Tucker 랭크 $r$의 함수로서
% PSNR과 압축률 $\rho$를 나타낸다.
% 높은 $r$은 더 나은 PSNR을 제공하지만 압축률은 낮아지는
% 기본적인 트레이드오프를 보여준다.

% \begin{figure}[h]
%   \centering
%   \fbox{\rule{0pt}{4cm}\rule{0.8\linewidth}{0pt}}  % 플레이스홀더
%   \caption{Tucker 랭크 $r$의 함수로서 PSNR (왼쪽 축)과 압축률 $\rho$ (오른쪽 축).
%            HOOI ($\mathcal{L}_F$, 실선)와 IRLS ($\mathcal{L}_1$, 점선).}
%   \label{fig:psnr_curve}
% \end{figure}

% \paragraph{시각적 비교.}
% 그림~\ref{fig:visual}은 원본 이미지와 $r \in \{8, 20, 40\}$에서의
% Tucker 재구성 결과, 그리고 오차 지도
% $|\tensor{X} - \hat{\tensor{X}}|$ (채널 평균)를 보여준다.

% \begin{figure}[h]
%   \centering
%   \fbox{\rule{0pt}{5cm}\rule{\linewidth}{0pt}}  % 플레이스홀더
%   \caption{원본 이미지 (왼쪽)와 랭크 $r = 8$, $r = 20$, $r = 40$에서의
%            Tucker 재구성 (2--4열). 아래 행: 점별 오차 지도 (hot 컬러맵).}
%   \label{fig:visual}
% \end{figure}

% \paragraph{손상 이미지에서의 노름 비교.}
% 그림~\ref{fig:norm_compare}은 10\% 희소 손상 하에서 HOOI, IRLS, Adam의
% 재구성 결과를 제시한다.
% 오차 지도는 \ref{subsec:norm_geometry}절에서 예측한 잔차 집중 효과를
% 부각시킨다: IRLS 오차는 손상 픽셀에 집중되는 반면,
% HOOI 오차는 균등하게 분산된다.

% \begin{figure}[h]
%   \centering
%   \fbox{\rule{0pt}{4cm}\rule{\linewidth}{0pt}}  % 플레이스홀더
%   \caption{10\% 희소 손상에서의 재구성.
%            왼쪽부터: 손상 입력, HOOI, IRLS, Adam. 아래 행: 오차 지도.}
%   \label{fig:norm_compare}
% \end{figure}

% \paragraph{특이값 스펙트럼.}
% 그림~\ref{fig:sv_spectrum}은 각 모드-$n$ 전개 $\unfold{X}{n}$,
% $n \in \{1, 2, 3\}$의 특이값 감소를 나타낸다.
% 감소율이 각 랭크 수준에서 포착되는 이미지 에너지의 양을 직접 결정하며,
% 오차 상한~\eqref{eq:hosvd_bound}을 검증한다.

% \begin{figure}[h]
%   \centering
%   \fbox{\rule{0pt}{4cm}\rule{0.7\linewidth}{0pt}}  % 플레이스홀더
%   \caption{테스트 이미지의 모드별 특이값 스펙트럼.
%            음영 영역 ($r < 20$)은 HOSVD 오차 상한에 기여하는 절단 특이값.}
%   \label{fig:sv_spectrum}
% \end{figure}

% % ============================================================
% \section{토의}
% \label{sec:discussion}
% % ============================================================

% \subsection{이론--실험 일치}
% \label{subsec:disc_theory}

% \paragraph{오차 상한 (실험 1).}
% 실험적 HOSVD 오차는 모든 테스트 랭크에서 이론적 상한
% $B(r,r,3)$~\eqref{eq:hosvd_bound} 아래로 일관되게 유지되었다.
% 특이값 스펙트럼이 천천히 감소하는 경우 상한이 타이트하지만,
% 급격한 스펙트럼 감소를 보이는 자연 이미지에서는
% 낮은 랭크에서 상한이 약 2--3배 느슨함을 확인하였다.
% HOOI는 모든 경우에서 HOSVD를 능가하며 단조 수렴을 확인하였다.

% \paragraph{노름 비교 (실험 2).}
% 결과는 표~\ref{tab:norm_noise}의 예측을 확인한다.
% 가우시안 잡음 하에서는 HOOI가 최저 RelErr를 달성하고 (MLE 최적성),
% 희소 손상 하에서는 IRLS가 손상 밀도 $p$에 따라 증가하는
% 격차로 HOOI를 능가한다.
% IRLS가 선호되기 시작하는 교차점 --- IRLS가 더 유리해지는 손상 밀도 ---은
% $p \approx 0.03$--0.05에서 관측되었으며, 이는 어떠한 $p > 0$에서도
% 프로베니우스 추정량이 이론적으로 분해된다는 예측과 일치한다.

% \paragraph{수렴 (실험 3).}
% HOOI는 모든 테스트 이미지에서 약 10--20회 반복으로 수렴하여
% \cite{de2000best}에 보고된 빠른 실용적 수렴을 확인하였다.
% Adam은 유사한 재구성 품질에 도달하기 위해 훨씬 더 많은 반복
% ($\sim$500회)이 필요하여, 닫힌 형태 ALS 갱신의 효율성 이점을 부각시킨다.
% HOOI로 웜 스타트한 IRLS는 15회 미만의 외부 반복으로 수렴하였다.

% \paragraph{초기화 (실험 4).}
% HOSVD 초기화는 무작위 초기화보다 일관되게 낮은 최종 오차를 달성하였으며,
% 분산이 없다 (결정론적). 무작위 직교 초기화는 무작위 가우시안보다
% 낮은 분산을 보여, 정확한 HOSVD 해를 사용할 수 없을 때도
% 직교성이 유용한 귀납적 편향임을 시사한다.

% \subsection{한계}
% \label{subsec:limitations}

% \begin{itemize}
%   \item \textbf{랭크 선택.}
%         Tucker 랭크 $(r_1, r_2, r_3)$은 특이값 감소에 기반하여
%         수동으로 선택하였다.
%         자동 랭크 선택 (예: 교차 검증, 정보 기준)은 다루지 않는다.

%   \item \textbf{CP 퇴화.}
%         CP-ALS의 알려진 수렴 어려움과 비고 2.3의 퇴화 문제로 인해
%         Tucker 분해에 집중하였으며 CP 분해는 구현하지 않았다.

%   \item \textbf{실험 규모.}
%         모든 실험은 단일 $256 \times 256$ 이미지를 사용한다.
%         다른 이미지, 크기, 도메인 (예: 초분광 영상, 비디오 텐서)으로의
%         일반화는 추가 연구가 필요하다.

%   \item \textbf{IRLS 수렴 보장.}
%         비볼록 설정에서 Tucker를 위한 IRLS의 외부 수렴은
%         이론적으로 보장되지 않으며, 관측된 수렴은 실험적이다.
% \end{itemize}

% % ============================================================
% \section{결론}
% \label{sec:conclusion}
% % ============================================================

% 본 보고서는 세 가지 재구성 노름 --- 프로베니우스, $\ell_1$, Huber ---과
% 두 가지 최적화 패러다임 --- ALS/HOOI 및 경사 기반 방법 ---
% 하에서 이미지 텐서의 저랭크 Tucker 근사를 연구하였다.

% 주요 발견은 다음과 같다.
% 첫째, HOSVD 오차 상한~\eqref{eq:hosvd_bound}은 재구성 오차의 신뢰할 수 있는
% 상한 추정값이며, HOOI는 이를 일관되게 개선한다.
% 둘째, 구조적 잡음 존재 시 재구성 노름의 선택이 중요하다:
% IRLS를 통한 $\ell_1$-Tucker는 희소 손상 하에서 현저히 우수한 재구성을
% 달성하는 반면, 프로베니우스-Tucker는 가우시안 잡음 하에서 최적이다.
% 셋째, HOSVD 초기화는 비볼록성의 영향을 줄이는 강력하고 결정론적인 출발점을
% 제공하며, 무작위 초기화는 문제 난이도에 따라 분산이 증가한다.
% 넷째, ALS/HOOI는 Adam에 비해 실시간 수렴 속도가 훨씬 빠르므로
% 프로베니우스 목적 함수에 선호되는 방법이며,
% 경사 기반 최적화는 닫힌 형태 ALS 갱신이 없는 비프로베니우스 손실에
% 필수적이다.

% 이러한 결과들은 분해 구조, 재구성 노름, 최적화 알고리즘 간의 상호작용이
% 실용적인 이미지 압축에 있어 비자명하고 중요함을 종합적으로 보여준다.
% 향후 연구로는 자동 랭크 선택, CP 및 텐서 트레인 분해로의 확장,
% 비디오 및 초분광 텐서 데이터 적용을 제안한다.